%----------------------------------------------------------------------------------------
%	PACKAGES AND OTHER DOCUMENT CONFIGURATIONS
%----------------------------------------------------------------------------------------

\documentclass[12pt]{article}

\usepackage{polski}
\usepackage[polish]{babel}
\usepackage[utf8]{inputenc}
\usepackage{datetime}
\usepackage{graphicx}
\usepackage{tikz}
\usepackage{amsmath}
\usepackage{amsfonts}
\usepackage{epstopdf}
\usepackage{multirow}
\usepackage{tabularx}
%\usepackage[colorlinks=true]{hyperref}
%\usepackage[all]{hypcap}
%\usepackage{showframe} 
\usepackage{geometry}
 \geometry{
 	a4paper, 
 	left	=20mm,
 	right	=20mm,
 	top		=20mm,
 	bottom	=20mm,
 }
 
%----------------------------------------------------------------------------------------
 
%----------------------------------------------------------------------------------------
% DATES
%----------------------------------------------------------------------------------------

\renewcommand{\dateseparator}{.}
\newdate{exercise_date}{20}{05}{2015}

%----------------------------------------------------------------------------------------

%----------------------------------------------------------------------------------------
% TIKZ PACKAGES
%----------------------------------------------------------------------------------------

\usetikzlibrary{arrows}

%----------------------------------------------------------------------------------------

\begin{document}
 
\begin{titlepage}

\newcommand{\HRule}{\rule{\linewidth}{0.5mm}}
% Defines a new command for the horizontal lines, change thickness here

\center
% Center everything on the page
 
%----------------------------------------------------------------------------------------
%	LOGO SECTION
%----------------------------------------------------------------------------------------

\includegraphics[width=6cm]{../res/img/logo.png}\\[1cm]
% Include a department/university logo - this will require the graphicx package
 
%----------------------------------------------------------------------------------------
 
%----------------------------------------------------------------------------------------
%	HEADING SECTIONS
%----------------------------------------------------------------------------------------

\textsc{\LARGE Akademia Górniczo-Hutnicza \\[0.2cm]
im. Stanisława Staszica w Krakowie}\\[1.5cm]
% Name of your university/college

\textsc{\Large Podstawy Automatyki}\\[0.5cm]
% Major heading such as course name

%----------------------------------------------------------------------------------------
%	TITLE SECTION
%----------------------------------------------------------------------------------------

\HRule \\[0.5cm]
{ \huge \bfseries Dyskretne układy regulacji \\[0.3cm] oraz \\[0.5cm] Analiza
serwomechanizmu \\[0.2cm] przekaźnikowego z wykorzystaniem płaszczyzny
fazowej}\\[0.3cm]
% Title of your document
\HRule \\[1.5cm]
 
%----------------------------------------------------------------------------------------
%	AUTHOR SECTION
%----------------------------------------------------------------------------------------

% \begin{minipage}{0.4\textwidth}
% \begin{flushleft} \large
% \emph{Author:}\\
% Konrad \textsc{Adasiewcz} % Your name
% \end{flushleft}
% \end{minipage}
% ~
% \begin{minipage}{0.4\textwidth}
% \begin{flushright} \large
% \emph{Supervisor:} \\
% dr inż. Paweł \textsc{Rotter} % Supervisor's Name
% \end{flushright}
% \end{minipage}\\[4cm]

% If you don't want a supervisor, uncomment the two lines below and remove the section above
\flushright
\Large \emph{Autorzy:}\\
Konrad \textsc{Adasiewcz}\\[0.1cm] % Your name
Michał \textsc{Maciejewski}\\[3cm] % Your name

%----------------------------------------------------------------------------------------
%	DATE SECTION
%----------------------------------------------------------------------------------------
Data wykonania ćwiczenia: \\
{\large \displaydate{exercise_date}}\\[1cm]


\vfill % Fill the rest of the page with whitespace

\end{titlepage}

\section*{Abstrakt}

Przed autorem został postawiony problem stworzenia aplikacji w pakiecie
\textsc{Matlab}, implementaującej algorytm największego spadku
do poszukiwania minimów, na przykładzie funkcji benchmarkowej Rosenbrocka.
Dodatkowym założeniem dla projektowanej aplikacji była maksymalna elastyczność w
postaci prostej podmiany minimalizowanej funkcji czy wyboru liczby wymiarów.

\section*{Wstęp}

Metoda największego spadku jest prostą iteracyjną metodą gradientową
poszukiwania minimów funkcjonałów określonych nad ciałem liczb rzeczywistych:

\begin{equation*}
    f: \mathbb{R}^{n} \rightarrow \mathbb{R} 
\end{equation*}

Dodatkowymi założeniami gwarantującymi poprawną zbieżność metody (gwarancja
zbieżności do minimum globalnego) są:

\begin{itemize}
  \item $f \in \mathcal{C}^{1}$ (funkcja ciągła i różniczkowalna)
  \item $f$ jest ściśle wypukła w badanym otoczeniu minimum globalnego
\end{itemize}

Jako że funkcja Rosenbrocka nie jest funkcją wypukłą na całej przestrzeni
$\mathbb{R}^{n}$ metoda nie gwarantuje zbieżności do minimum globalnego (dla
pewnych punktów startowych algorytm może zbiegać do minimum lokalnego, lub
punktu przegięcia funkcjonału).

\section*{Algorytm metody największego spadku}

Badana metoda jest metodą gradientową, iteracyjną, co oznacza że kolejne
przybliżenie minimum funkcji jest wyznaczane krokowo na podstawie poprzedniego
punktu oraz wartości gradientu w danym punkcie, co można przedstawić za pomocą
rekurencyjnego równania:

\begin{equation}
    x_{k+1} = \min_{\alpha \in [0, \alpha_{\textrm{m}}]} \left[x_{k} - \alpha \nabla f(x_{k})\right]
\end{equation}

Zatem problem minimalizacji funkcjonału w przestrzeni wielowymiarowej zostaje
sprowadzony do wielokrotnej minimalizacji funkcji jednej zmiennej $\alpha$.


\end{document}